\documentclass[12pt,a4paper,oneside]{article}

\usepackage[utf8]{inputenc}
\usepackage[T1]{fontenc}
\usepackage[brazil]{babel}
\usepackage{geometry}
\geometry{a4paper, margin=2.5cm}
\usepackage{indentfirst}
\usepackage{times}
\usepackage{hyperref}
\usepackage{booktabs}
\usepackage{array}
\usepackage{graphicx}
\usepackage{abntex2cite}

\setlength{\parindent}{1.25cm}
\setlength{\parskip}{0.3cm}

\begin{document}

\begin{center}
  \textbf{\large ESCOLA E FACULDADE DE TECNOLOGIA SENAI GASPAR RICARDO JÚNIOR}\\
  \textbf{Desenvolvimento de Sistemas}\\[1cm]
  \textbf{\Large Lucas Miguel Leite}\\[0.5cm]
  \textbf{Professor: Deivison Takatu}\\[2cm]
  \textbf{\LARGE\textbf{Sistemas ERP na Integração Industrial}}\\[1cm]
  Sorocaba -- 06/03/2026
\end{center}

\vspace{1cm}

\begin{center}
  \textbf{RESUMO}
\end{center}

Este artigo analisa o papel dos sistemas ERP (\textit{Enterprise Resource Planning}) como integradores estratégicos na indústria, posicionando-os no nível mais alto da Pirâmide da Automação. O estudo aborda a definição e as características principais dos ERPs, a estrutura de módulos (financeiro, compras, estoque, produção, vendas e recursos humanos) e, sobretudo, a integração automática entre esses módulos. Através de uma tabela detalhada de eventos e impactos cruzados, demonstra-se como um único evento em um setor gera efeitos em cascata por toda a organização. O exemplo prático de uma fábrica de bebidas ilustra o fluxo completo desde o pedido do cliente até o retorno dos indicadores de produção. Conclui-se que o ERP funciona como o sistema nervoso da organização, consolidando a integração vertical ao conectar operação e estratégia de forma contínua.

\vspace{0.5cm}
\textbf{Palavras-chave:} ERP. Integração de Sistemas. Gestão Empresarial. Módulos ERP. PCP.

\newpage

\section{Introdução}

Na evolução da disciplina de Integração Vertical e Horizontal, compreendeu-se progressivamente o papel de cada nível da automação: redes industriais para comunicação, SDCD/CLP para controle, SCADA para supervisão e MES para execução. Chega-se agora ao nível mais alto da pirâmide: o ERP (\textit{Enterprise Resource Planning}) \cite{caiçara2015}. O ERP é responsável por integrar e gerenciar toda a empresa --- produção, logística, financeiro, compras, vendas e recursos humanos --- transformando dados operacionais em decisões estratégicas.

\section{Conceito de ERP}

\subsection{Definição}

O ERP é um sistema integrado de gestão que centraliza informações e processos de toda a organização em uma única plataforma. Possui base de dados única, módulos integrados, atualização em tempo real e padronização de processos, o que evita retrabalho, inconsistência de dados e falhas de comunicação entre setores \cite{santos2024}.

\subsection{ERP como Sistema Integrador}

Sem ERP, os setores operam de forma isolada, com planilhas desconectadas e decisões imprecisas. Com ERP, a empresa adquire visão global, dados consistentes e capacidade de resposta rápida. O ERP atua no nível estratégico e gerencial, respondendo a perguntas fundamentais: o que produzir? Quando? Quanto? A que custo? \cite{groover2011}.

\section{Módulos de um ERP}

\subsection{Principais Módulos}

\begin{itemize}
  \item \textbf{Financeiro:} Contas a pagar/receber, fluxo de caixa, faturamento.
  \item \textbf{Compras:} Gestão de fornecedores, pedidos de compra, controle de suprimentos.
  \item \textbf{Estoque:} Controle de materiais, entrada e saída, inventário.
  \item \textbf{Produção (PCP):} Planejamento da produção, ordens de produção, controle de capacidade.
  \item \textbf{Vendas:} Pedidos de clientes, faturamento, histórico de vendas.
  \item \textbf{Recursos Humanos:} Folha de pagamento, gestão de colaboradores, avaliação de desempenho.
\end{itemize}

\subsection{Integração entre Módulos}

A principal força de um ERP está na integração automática entre módulos, permitindo que uma ação em um setor gere impactos imediatos em toda a empresa. A Tabela~\ref{tab:integracao} demonstra essa dinâmica.

\begin{table}[htb]
  \centering
  \caption{Integração entre módulos do ERP: eventos e impactos cruzados.}
  \label{tab:integracao}
  \resizebox{\textwidth}{!}{%
  \begin{tabular}{p{4cm}p{3cm}p{4cm}p{4.5cm}}
    \toprule
    \textbf{Evento} & \textbf{Módulo de Origem} & \textbf{Módulos Impactados} & \textbf{Resultado} \\
    \midrule
    Venda realizada & Vendas & Financeiro, Estoque & Geração de receita e baixa no estoque \\
    Produção iniciada & Produção (PCP) & Estoque & Consumo de matéria-prima \\
    Produção finalizada & Produção & Estoque, Vendas & Entrada de produto acabado \\
    Estoque mínimo atingido & Estoque & Compras & Geração de pedido de compra \\
    Pedido de compra aprovado & Compras & Financeiro & Contas a pagar \\
    Alocação de operador & RH & Produção & Atualização de custo operacional \\
    Atraso na produção & Produção & Vendas, Logística & Reprogramação de entrega \\
    Devolução de produto & Vendas & Estoque, Financeiro & Ajuste de estoque e estorno \\
    \bottomrule
  \end{tabular}%
  }
\end{table}

Nenhum módulo funciona isoladamente: um único evento pode impactar vários setores simultaneamente, e a informação flui automaticamente pela empresa \cite{cardoso2021}.

\section{Exemplo Prático: Fábrica de Bebidas}

O fluxo completo ilustra a integração vertical proporcionada pelo ERP:

\begin{enumerate}
  \item \textbf{Pedido:} Cliente solicita 10.000 unidades. O ERP registra a venda, calcula o custo e verifica o estoque.
  \item \textbf{Planejamento:} O ERP gera a ordem de produção e a envia ao MES.
  \item \textbf{Execução:} O MES distribui a produção e planeja a utilização das máquinas.
  \item \textbf{Produção:} CLPs e sensores executam o processo industrial.
  \item \textbf{Monitoramento:} O SCADA exibe dados em tempo real.
  \item \textbf{Retorno:} O ERP recebe dados de produção realizada, custos, tempo e outros indicadores.
\end{enumerate}

\section{Conclusão}

O ERP consolida a integração da empresa ao transformar dados operacionais em informação gerencial útil. Seu principal valor reside na capacidade de alinhar planejamento e execução, dar visibilidade do negócio como um todo e permitir respostas rápidas a mudanças. Na prática, o ERP garante que o que foi planejado realmente aconteça --- e que possa ser ajustado com base em dados reais. Assim, o ERP fecha o ciclo da integração vertical, conectando operação e estratégia de forma contínua \cite{schwab2016}.

\bibliographystyle{plain}
\bibliography{referencias}

\end{document}
